%% ma: L11-B21-0015
%% lop: 11
%% chuong: 6
%% bai: 21
%% dang: 11.21.10
%% loai: TLN
%% muc: VD
%% dap_an: "30"
%% nguon: "(NBV)-ĐỀ SỐ 15-MỖI NGÀY 1 ĐỀ THI 2025.pdf (D15, MỖI NGÀY 1 ĐỀ - Đề số 15), Phần III Câu 1"
%% kien_thuc: "Bài 7 (cấp số nhân), Bài 21 (bất phương trình mũ, lôgarit)"
%% trang_thai: cho-duyet
%% ly_do: "Dạng mới 11.21.10 đề xuất (lãi kép theo kì hạn: tìm thời gian gửi để đạt số tiền mong muốn bằng bất phương trình mũ)."
%% ghi_chu: "Đáp án gốc 30 đúng. Bỏ ảnh minh họa quầy ngân hàng (ảnh trang trí)."
\begin{ex}
Một người gửi tiết kiệm $100$ triệu đồng vào ngân hàng theo thể thức lãi kép kì hạn $6$ tháng với lãi suất $8\%$ một năm. Giả sử lãi suất không thay đổi. Hỏi sau bao nhiêu tháng người đó nhận được ít nhất $120$ triệu đồng gồm cả gốc và lãi?
\shortans{30}
\loigiai{Lãi suất mỗi kì $6$ tháng là $\dfrac{8\%}{2}=4\%$. Sau $n$ kì, số tiền là $100\cdot(1{,}04)^n$ (triệu đồng).

$100\cdot(1{,}04)^n\ge120\iff n\ge\log_{1{,}04}1{,}2\approx4{,}65$, nên $n=5$ kì, tức $5\cdot6=30$ tháng.}
\end{ex}
